% =====================================================================
% Apêndice B — de onde vem cada número
% =====================================================================

\chapter{Proveniência dos resultados preliminares}
\label{ap:proveniencia}

Os resultados do Capítulo~\ref{cap:resultados} vêm de uma única execução do
pipeline com dados reais, em 21 e 22 de setembro de 2026, na máquina do
pesquisador. Os microdados e os arquivos intermediários não são versionados. Os
números foram transcritos para o texto do Ensaio 1 e, dele, para este projeto;
nenhum número novo foi produzido para este documento. O
Quadro~\ref{qua:proveniencia} indica o programa que gera cada resultado, para
que a nova execução prevista no plano de trabalho possa ser conferida contra
esta.

\begin{quadro}[htb]
\centering
\caption{Programa de origem de cada resultado}
\label{qua:proveniencia}
\small
\begin{tabular}{|p{6.4cm}|p{7.8cm}|}
\hline
\textbf{Resultado} & \textbf{Programa (pasta \texttt{scripts/})} \\
\hline
Nível e curva dose-resposta & \texttt{estimate/03\_contdid.R} \\
\hline
Inclinação em dose e três procedimentos de inferência &
  \texttt{estimate/04\_robustness.py} \\
\hline
Poder do teste de tendências anteriores & \texttt{estimate/06\_pretrends.R} \\
\hline
Sensibilidade a violações de paralelismo & \texttt{estimate/07\_honestdid.R} \\
\hline
Controle sintético & \texttt{estimate/08\_synthdid.R} \\
\hline
Correção de Holm & \texttt{estimate/09\_holm.py} \\
\hline
Inclinação em datas fictícias & \texttt{estimate/10\_spt\_pretrend.py} \\
\hline
Erro de classificação e calibração &
  \texttt{estimate/12\_erro\_de\_classificacao.py} \\
\hline
Custos da Conab & \texttt{data\_prep/12\_clean\_conab\_custos.py} \\
\hline
Efeito mínimo detectável do desenho de calendário (a executar) &
  \texttt{estimate/14\_mde\_calendario.py} \\
\hline
\end{tabular}
\legend{Fonte: elaboração própria.}
\end{quadro}
